\chapter{Технологическая часть}

\section{Средства разработки}

В качестве языка программирования со статической типизацией был выбран язык C++ (стандарт C++20)~\cite{iso_cpp}. В процессе разработки использовались следующие средства стандартной библиотеки языка:
\begin{itemize}
    \item контейнер \texttt{std::vector<std::vector<int>>} --- для динамического представления целочисленных матриц;
    \item контейнер \texttt{std::vector<int>} --- для хранения вспомогательных массивов $row$ и $col$;
    \item функция \texttt{clock()} из заголовочного файла \texttt{<ctime>} --- для замера процессорного времени работы исследуемых алгоритмов;
\end{itemize}

Сборка программы осуществлялась компилятором GCC с флагами компиляции:
\begin{center}
\texttt{-std=c++20 -Wall -Wextra -Wpedantic -O0}
\end{center}
Флаг оптимизации \texttt{-O0} задан принудительно, чтобы компилятор не выполнял автоматическую оптимизацию.

\section{Реализация алгоритмов}

На рисунке~\ref{fig:std_mult} представлена программная реализация стандартного алгоритма умножения матриц.
\begin{figure}[H]
    \centering
    \lstinputlisting[
        firstline=5, lastline=19
    ]{multiply.cpp}
    \caption{Стандартный алгоритм умножения матриц}
    \label{fig:std_mult}
\end{figure}

На рисунке~\ref{fig:winograd_mult} представлена реализация базового алгоритма Винограда.
\begin{figure}[H]
    \centering
    \lstinputlisting[
        firstline=21, lastline=54
    ]{multiply.cpp}
    \caption{Базовый алгоритм Винограда}
    \label{fig:winograd_mult}
\end{figure}

На рисунке~\ref{fig:winograd_opt_mult} представлена реализация оптимизированного алгоритма Винограда.
\begin{figure}[H]
    \centering
    \lstinputlisting[
        firstline=56, lastline=103
    ]{multiply.cpp}
    \caption{Оптимизированный алгоритм Винограда}
    \label{fig:winograd_opt_mult}
\end{figure}

\section{Тестирование реализаций}

В таблице~\ref{tbl:tests} представлены результаты функционального тестирования разработанных реализаций алгоритмов умножения матриц.

\begingroup
\setlength{\arraycolsep}{2.5pt}

\begin{table}[H]
    \caption{Функциональные тесты}
    \label{tbl:tests}
    \centering
    \begin{tblr}{
            width = \linewidth,
            colsep = 3pt,
            colspec = {Q[c,m] Q[c,m] Q[c,m] Q[c,m] Q[c,m] X[c,m]},
            hlines, vlines,
            row{1,2} = {font=\bfseries\small},
            row{3-Z} = {font=\small},
        }
        \SetCell[r=2]{c} №
        & \SetCell[c=2]{c} Входные данные &
        & \SetCell[c=2]{c} Выходные данные &
        & \SetCell[r=2]{c} Комментарий \\
        & Матрица $A$ & Матрица $B$ & Ожид. & Получ. & \\
        1  & $\begin{pmatrix} 3 \end{pmatrix}$
           & $\begin{pmatrix} -4 \end{pmatrix}$
           & $\begin{pmatrix} -12 \end{pmatrix}$
           & $\begin{pmatrix} -12 \end{pmatrix}$
           & Матрицы размера $1 \times 1$ \\
        2  & $\begin{pmatrix} 1 & 2 & 3 \\ 4 & 5 & 6 \\ 7 & 8 & 9 \end{pmatrix}$
           & $\begin{pmatrix} 1 & 0 & 0 \\ 0 & 1 & 0 \\ 0 & 0 & 1 \end{pmatrix}$
           & $\begin{pmatrix} 1 & 2 & 3 \\ 4 & 5 & 6 \\ 7 & 8 & 9 \end{pmatrix}$
           & $\begin{pmatrix} 1 & 2 & 3 \\ 4 & 5 & 6 \\ 7 & 8 & 9 \end{pmatrix}$
           & Умножение на единичную матрицу $E_{3 \times 3}$ \\
        3  & $\begin{pmatrix} 1 & 2 \\ 3 & 4 \end{pmatrix}$
           & $\begin{pmatrix} 0 & 0 \\ 0 & 0 \end{pmatrix}$
           & $\begin{pmatrix} 0 & 0 \\ 0 & 0 \end{pmatrix}$
           & $\begin{pmatrix} 0 & 0 \\ 0 & 0 \end{pmatrix}$
           & Умножение на нулевую матрицу \\
        4  & $\begin{pmatrix} 1 & 2 \\ 3 & 4 \end{pmatrix}$
           & $\begin{pmatrix} 5 & 6 \\ 7 & 8 \end{pmatrix}$
           & $\begin{pmatrix} 19 & 22 \\ 43 & 50 \end{pmatrix}$
           & $\begin{pmatrix} 19 & 22 \\ 43 & 50 \end{pmatrix}$
           & Квадратные матрицы $2 \times 2$, чётный размер \\
        5  & $\begin{pmatrix} 1 & 2 & 3 \\ 4 & 5 & 6 \\ 7 & 8 & 9 \end{pmatrix}$
           & $\begin{pmatrix} 9 & 8 & 7 \\ 6 & 5 & 4 \\ 3 & 2 & 1 \end{pmatrix}$
           & $\begin{pmatrix} 30 & 24 & 18 \\ 84 & 69 & 54 \\ 138 & 114 & 90 \end{pmatrix}$
           & $\begin{pmatrix} 30 & 24 & 18 \\ 84 & 69 & 54 \\ 138 & 114 & 90 \end{pmatrix}$
           & Квадратные матрицы $3 \times 3$, нечётный размер \\
        6  & $\begin{pmatrix} 1 & 2 & 3 \\ 4 & 5 & 6 \end{pmatrix}$
           & $\begin{pmatrix} 7 & 8 \\ 9 & 10 \\ 11 & 12 \end{pmatrix}$
           & $\begin{pmatrix} 58 & 64 \\ 139 & 154 \end{pmatrix}$
           & $\begin{pmatrix} 58 & 64 \\ 139 & 154 \end{pmatrix}$
           & Прямоугольные матрицы $2 \times 3$ и $3 \times 2$ \\
        7  & $\begin{pmatrix} 1 & 2 \\ 3 & 4 \\ 5 & 6 \end{pmatrix}$
           & $\begin{pmatrix} 1 & 2 & 3 & 4 \\ 5 & 6 & 7 & 8 \end{pmatrix}$
           & $\begin{pmatrix} 11 & 14 & 17 & 20 \\ 23 & 30 & 37 & 44 \\ 35 & 46 & 57 & 68 \end{pmatrix}$
           & $\begin{pmatrix} 11 & 14 & 17 & 20 \\ 23 & 30 & 37 & 44 \\ 35 & 46 & 57 & 68 \end{pmatrix}$
           & Прямоугольные матрицы $3 \times 2$ и $2 \times 4$ \\
        8  & $\begin{pmatrix} 1 & 2 & 3 & 4 \end{pmatrix}$
           & $\begin{pmatrix} 1 \\ 2 \\ 3 \\ 4 \end{pmatrix}$
           & $\begin{pmatrix} 30 \end{pmatrix}$
           & $\begin{pmatrix} 30 \end{pmatrix}$
           & Вектор-строка на вектор-столбец \\
        9  & $\begin{pmatrix} 1 \\ 2 \\ 3 \end{pmatrix}$
           & $\begin{pmatrix} 4 & 5 \end{pmatrix}$
           & $\begin{pmatrix} 4 & 5 \\ 8 & 10 \\ 12 & 15 \end{pmatrix}$
           & $\begin{pmatrix} 4 & 5 \\ 8 & 10 \\ 12 & 15 \end{pmatrix}$
           & Вектор-столбец на вектор-строку \\
    \end{tblr}
\end{table}

\endgroup

Все тесты были пройдены успешно без расхождений в результатах.

\section{Вывод}

В результате выполнения технологической части были обоснованы средства реализации, приведены листинги исходного кода стандартного, базового и оптимизированного алгоритмов Винограда на языке C++, а также разработана система тестов, подтвердившая корректность всех разработанных реализаций.